\chapter{พจนานุกรมคำศัพท์เทคนิคและนิยามศัพท์}
\label{app:glossary}
\index{พจนานุกรมคำศัพท์}
\index{นิยามศัพท์}

ภาคผนวกนี้รวบรวมคำศัพท์เทคนิคเฉพาะทางด้านคอมพิวเตอร์วิทัศน์ การเรียนรู้เชิงลึก พยาธิวิทยาพืช และเกษตรแม่นยำ พร้อมคำนิยามและความหมายภาษาไทย

\begin{description}[leftmargin=3.5cm, style=nextline]
    \item[Backpropagation (การแพร่ย้อนกลับ)] กระบวนการคำนวณอนุพันธ์ย่อยของฟังก์ชันการสูญเสียเทียบกับค่าน้ำหนักทุกตัวในโครงข่ายประสาทเทียม เพื่อใช้ในการปรับปรุงค่าน้ำหนักผ่านขั้นตอนวิธี Gradient Descent
    \item[Batch Normalization (การปรับมาตรฐานบัตช์)] เทคนิคการปรับค่าข้อมูลในมินิบัตช์ให้มีค่าเฉลี่ยเป็นศูนย์และส่วนเบี่ยงเบนมาตรฐานเป็นหนึ่ง เพื่อช่วยลดปัญหา Internal Covariate Shift และทำให้โมเดลลู่เข้าเร็วขึ้น
    \item[Bounding Box (กรอบสี่เหลี่ยมระบุพิกัด)] พิกัดสี่เหลี่ยม $(x_{\text{min}}, y_{\text{min}}, x_{\text{max}}, y_{\text{max}})$ ที่ใช้ระบุตำแหน่งและขอบเขตของวัตถุหรือรอยโรคบนภาพดิจิทัล
    \item[Confusion Matrix (เมทริกซ์ความสับสน)] ตารางสรุปผลการทำนายของโมเดลเทียบกับค่าความจริง ซึ่งประกอบด้วยช่อง True Positive, False Positive, True Negative และ False Negative
    \item[Depthwise Separable Convolution (การแยกมิติคอนโวลูชันเชิงลึก)] เทคนิคการแยกการคำนวณคอนโวลูชันออกเป็น Depthwise Conv (กรองเชิงพื้นที่) และ Pointwise Conv (รวมช่องสัญญาณ) เพื่อลดปริมาณการคำนวณลง 8--9 เท่า
    \item[Early Stopping (การหยุดก่อนกำหนด)] เทคนิค Regularization ที่ติดตามค่า Validation Loss และยุติการฝึกฝนทันทีเมื่อผลการประเมินเริ่มแย่ลงต่อเนื่องตามจำนวนรอบที่กำหนด เพื่อป้องกัน Overfitting
    \item[Edge Computing (การประมวลผลบนขอบ)] สถาปัตยกรรมการประมวลผลข้อมูลที่ดำเนินการบนอุปกรณ์ปลายทางโดยตรง เช่น สมาร์ทโฟน หรือบอร์ดสมองกลฝังตัว โดยไม่ต้องพึ่งพาเซิร์ฟเวอร์คลาวด์
    \item[Explainable AI - XAI (ปัญญาประดิษฐ์อธิบายได้)] กลุ่มเทคนิคทางปัญญาประดิษฐ์ที่มุ่งเน้นการเปิดเผยเหตุผล กลไก หรือพื้นที่ที่โมเดลใช้ในการตัดสินใจ เพื่อสร้างความโปร่งใสและความน่าเชื่อถือให้แก่มนุษย์
    \item[Focal Loss (ฟังก์ชันการสูญเสียแบบโฟคอล)] ฟังก์ชัน Loss ที่ปรับปรุงจาก Cross-Entropy โดยเพิ่มตัวคูณถ่วงน้ำหนักเพื่อลดความสำคัญของตัวอย่างที่จำแนกได้ง่าย และเพิ่มความสำคัญให้กับตัวอย่างที่จำแนกได้ยาก
    \item[Grad-CAM (การสร้างแผนที่ความร้อนถ่วงน้ำหนักเกรเดียนต์)] เทคนิค XAI ที่ใช้เกรเดียนต์ของคลาสเป้าหมายเทียบกับเลเยอร์คอนโวลูชันสุดท้าย เพื่อสร้างภาพความร้อนระบุบริเวณที่ส่งผลต่อการตัดสินใจ
    \item[Inverted Residuals (โครงสร้างตกค้างแบบกลับด้าน)] บล็อกโครงสร้างใน MobileNetV2/V3 ที่ทำการขยายช่องสัญญาณก่อนทำการคอนโวลูชันเชิงลึก จากนั้นจึงบีบอัดกลับลงมา เพื่อรักษาข้อมูลไม่ให้สูญหาย
    \item[Mycelium (เส้นใยรา)] โครงสร้างเส้นสายของเชื้อราที่เจริญแทรกซึมเข้าสู่เนื้อเยื่อพืชเพื่อดูดซับสารอาหารและปล่อยเอนไซม์ย่อยสลายเซลล์
    \item[ONNX - Open Neural Network Exchange] มาตรฐานเปิดสำหรับแลกเปลี่ยนและจัดเก็บโมเดลโครงข่ายประสาทเทียม ช่วยให้สามารถถ่ายโอนโมเดลข้ามเฟรมเวิร์กและรันบนเอนจินต่างๆ ได้อย่างอิสระ
    \item[Quantization (การลดทอนความแม่นยำตัวเลข)] กระบวนการแปลงค่าพารามิเตอร์ของโมเดลจากจำนวนจริง Floating-point (FP32) ไปเป็นจำนวนเต็ม Integer (INT8) เพื่อลดขนาดไฟล์และเพิ่มความเร็วในการประมวลผล
    \item[Squeeze-and-Excitation - SE Module] กลไกความใส่ใจที่คำนวณและปรับค่าน้ำหนักความสำคัญของแต่ละช่องสัญญาณ (Channel-wise Attention) โดยอัตโนมัติ
    \item[Transfer Learning (การเรียนรู้แบบถ่ายโอน)] กระบวนการนำโครงข่ายประสาทเทียมที่ผ่านการฝึกฝนบนชุดข้อมูลขนาดใหญ่มาปรับใช้กับงานใหม่ที่มีข้อมูลจำกัด โดยคงค่าน้ำหนักตัวสกัดคุณลักษณะเดิมไว้
\end{description}
